\chapter{Proposiciones}

  Existen 2 tipos de proposiciones, atómicas y moleculares. Estas deben ser
  fórmulas bien formadas, es decir, solo pueden tener un valor de verdad
  asociado, y su enunciado no puede contener valores de verdad.

  \section{Conectores}

  En lógica, se definen conectores; estos representan un par ordenado de
  valores asociados a la relación entre 2 proposiciones atómicas. Estas
  definiciones indican qué casos deben cumplirse del enunciado, y estos
  pares ordenados se pueden representar por valores en una tabla de verdad.

  \subsection{Conjunción}

  \begin{minipage}[t]{0.48\textwidth}
  \begin{definition}[title={Conjunción}]
  El conector $p \land q$ solo permite el par ordenado $(V, V)$ y excluye
  al resto.
  \end{definition}
  \end{minipage}

  \begin{minipage}[t]{0.48\textwidth}
  \begin{tabular}{c|c|c}
  $p$ & $q$ & $p \land q$ \\
  \hline
  V & V & V \\
  V & F & F \\
  F & V & F \\
  F & F & F \\
  \end{tabular}
  \end{minipage}

  De este modo, podemos probar la conjunción como verdad por una prueba
  directa. Para probar que es falsa, basta probar un contraejemplo:

  \begin{align}
  \sim p &\to \sim(p \land q) \\
  \sim q &\to \sim(p \land q)
  \end{align}

  \subsection{Disyunción}

  \begin{minipage}[t]{0.48\textwidth}
  \begin{definition}[title={Disyunción}]
  Para el conector $p \lor q$ únicamente se excluye $(F, F)$.
  \end{definition}
  \end{minipage}

  \begin{minipage}[t]{0.48\textwidth}
  \begin{tabular}{c|c|c}
  $p$ & $q$ & $p \lor q$ \\
  \hline
  V & V & V \\
  V & F & V \\
  F & V & V \\
  F & F & F \\
  \end{tabular}
  \end{minipage}

  Dado que la lógica no funciona por enumeración, la forma correcta es
  por exclusión del contraejemplo, así basta probar que ocurre:

  \begin{equation}
  \sim p \land \sim q \equiv (F, F)
  \end{equation}

  \subsection{Implicación}

  \begin{minipage}[t]{0.48\textwidth}
  \begin{definition}[title={Implicación}]
  En lógica proposicional no se busca describir el mundo físico, es decir,
  la lógica no entiende causalidad ni temporalidad, sino que son un
  conjunto de reglas. Este conector no explica que si ocurre $p$ causa $q$,
  no es que para que ocurra $q$ deba ocurrir $p$. No se evalúa de forma
  secuencial, sino que cada uno tiene su valor. La única exclusión que
  tiene es el caso $(V, F)$.
  \end{definition}
  \end{minipage}

  \begin{minipage}[t]{0.48\textwidth}
  \begin{tabular}{c|c|c}
  $p$ & $q$ & $p \to q$ \\
  \hline
  V & V & V \\
  V & F & F \\
  F & V & V \\
  F & F & V \\
  \end{tabular}
  \end{minipage}

  Dicho de otra forma, si $p$ es falso, significa que la hipótesis no se
  cumple, por tanto no hace falta evaluar el valor de $q$. La regla de
  exclusión dice:

  \alert{No existe un caso en que $p$ sea $V$ y $q$ sea $F$}

  \alert{Si no hay promesa, no hay incumplimiento}

  El caso $(F, F)$ puede parecer contraintuitivo; ambas falsas debería ser
  falsa la proposición en general, pero esto no es así porque no se busca
  reflejar o explicar una situación del mundo físico. Esta proposición
  también se puede abordar como:

  \begin{equation}
  \sim(p \land \sim q) \equiv p \to q
  \end{equation}

  Para demostrar la identidad de dos proposiciones basta con ver sus
  valores de verdad.

  \begin{tabular}{c|c|c|c|c|c}
  $p$ & $q$ & $p \to q$ & $\sim q$ & $p \land \sim q$ & $\sim(p \land \sim q)$ \\
  \hline
  V & V & V & F & F & V \\
  V & F & F & V & V & F \\
  F & V & V & F & F & V \\
  F & F & V & V & F & V \\
  \end{tabular}

  \section{Contraejemplo}

  Estas definiciones nos permiten identificar algunas estrategias para
  demostrar proposiciones matemáticas. Para ello, basta probar, ya fuese
  un \textit{contraejemplo} o una construcción directa.

  \begin{definition}[title={Contraejemplo}]
  Un contraejemplo es probar que alguna de las exclusiones puede ocurrir.
  \end{definition}
